% ============================================================
%  Póster 70 × 100 cm (vertical) — ¿Cómo estimar el índice de Theil con la GEIH?
%  Proyecto de Estadística (73210011), Universidad del Rosario, 2026-II
%
%  Las tablas y figuras vienen de outputs/ (las genera R con run_all.R).
%  Compilar desde esta carpeta: pdflatex -> biber -> pdflatex
%  Regla práctica: poco texto, letra grande, una figura por bloque.
% ============================================================
\documentclass[final]{beamer}
\usepackage[orientation=portrait, size=custom, width=70, height=100, scale=1.4]{beamerposter}

\usepackage[T1]{fontenc}
\usepackage[utf8]{inputenc}
\usepackage[spanish,es-nodecimaldot,es-noshorthands]{babel}
\usepackage{newtxtext,newtxmath}
\usepackage{icomma}
\usepackage{booktabs}

\setbeamertemplate{navigation symbols}{}
\setbeamercolor{block title}{fg=white, bg=red!60!black}
\setbeamercolor{block body}{bg=white}

\graphicspath{{../outputs/figuras/}{../docs/curso/report/}}
% Figura con alto máximo para que todo quepa en los 100 cm
\newcommand{\figura}[1]{\includegraphics[width=\linewidth,height=0.15\paperheight,keepaspectratio]{#1}}
\newcommand{\tabla}[1]{\input{../outputs/tablas/#1.tex}}
\newcommand{\tablaancha}[1]{\resizebox{\linewidth}{!}{\tabla{#1}}}

\usepackage{csquotes}
\usepackage[backend=biber,style=apa]{biblatex}
\addbibresource{../report/references.bib}

\begin{document}
\begin{frame}[t]

  % ---------------- Encabezado ----------------
  \begin{columns}[c]
    \begin{column}{0.18\paperwidth}
      \includegraphics[width=\linewidth]{ur.png}
    \end{column}
    \begin{column}{0.72\paperwidth}
      \centering
      {\Huge\bfseries ¿Cómo estimar el índice de Theil con la GEIH?\par}
      \vspace{0.5ex}
      {\LARGE Desigualdad del ingreso laboral en Colombia, 2018--2026\par}
      \vspace{0.5ex}
      {\Large Gabriel Salcedo · Jerónimo Marín · Isaac Suarez · Santiago Velandia\par}
      {\large Universidad del Rosario --- Estadística (73210011)\par}
    \end{column}
  \end{columns}

  \vspace{1cm}

  % ---------------- Cuerpo en dos columnas ----------------
  \begin{columns}[t]

    \begin{column}{0.45\paperwidth}

      \begin{block}{Problema y pregunta}
        % Dos o tres frases y la pregunta de investigación.
      \end{block}

      \begin{block}{Datos y método}
        % GEIH, población, factor de expansión, los cuatro estimadores (A--D),
        % diseño de la simulación.
      \end{block}

      \begin{block}{¿Cómo es el ingreso laboral?}
        \centering
        \figura{07a_histograma_ingreso}
      \end{block}

      \begin{block}{¿Es lognormal?}
        \centering
        \figura{p7_ln_ingreso_normalidad}
      \end{block}

    \end{column}

    \begin{column}{0.45\paperwidth}

      \begin{block}{Simulación Monte Carlo (H1)}
        % Figura de ECM frente a n. PENDIENTE: la simulación todavía no existe.
      \end{block}

      \begin{block}{Desigualdad 2018--2026 (H3)}
        \centering
        \tabla{theil_por_anio}
      \end{block}

      \begin{block}{¿Qué explica la desigualdad? (H2)}
        \centering
        \tabla{theil_participacion_entre}
      \end{block}

      \begin{block}{Inferencia clásica}
        \tablaancha{resumen_pruebas}
      \end{block}

      \begin{block}{Conclusiones}
        % Tres viñetas.
      \end{block}

      \begin{block}{Referencias}
        % Solo aparecen las obras citadas con \textcite{} o \parencite{} en el póster.
        % Elegir 3 o 4: no cabe la lista completa.
        \small
        \printbibliography[heading=none]
      \end{block}

    \end{column}

  \end{columns}

\end{frame}
\end{document}
